% ============================================================
% Section 72: 铜的费米能与导带宽度
% ============================================================
\section{固体理论：铜的费米能与导带宽度}
\label{sec:copper-fermi-energy}

\begin{modelbox}[题目：单价金属铜的自由电子参数计算]
铜是单价金属，密度为 $\rho=8000\,\mathrm{kg/m^3}$，原子量为 $64$。
\begin{enumerate}
    \item 计算绝对零度时电子的费米能；
    \item 估算导带的宽度。
\end{enumerate}
\end{modelbox}

\begin{tipbox}[考点快速分析]
\begin{itemize}
    \item \textbf{晶体结构}：铜为 FCC，每个惯用晶胞含 4 个原子
    \item \textbf{电子密度}：由宏观密度和原子量求晶格常数，再求电子数密度
    \item \textbf{费米能}：$E_F = \hbar^2k_F^2/(2m)$，其中 $k_F=(3\pi^2n)^{1/3}$
    \item \textbf{导带宽度}：自由电子模型下，取布里渊区边界最远点的能量
\end{itemize}
\end{tipbox}

\begin{derivationbox}[标准解题过程]
\textbf{Step 1: 电子数密度与晶格常数}

铜为面心立方（FCC）结构，每个惯用晶胞含 4 个铜原子。设晶格常数为 $a$，则密度与原子质量的关系为
\begin{equation}
\rho = \frac{4\times(M/N_A)}{a^3},
\end{equation}
其中 $M=64\,\mathrm{g/mol}=0.064\,\mathrm{kg/mol}$ 为摩尔质量，$N_A=6.02\times10^{23}\,\mathrm{mol^{-1}}$。

解出晶格常数：
\begin{align}
a^3 &= \frac{4M}{\rho N_A}
= \frac{4\times0.064}{8000\times6.02\times10^{23}}
= \frac{0.256}{4.816\times10^{27}}
\approx 5.32\times10^{-29}\,\mathrm{m^3}, \notag\\
a &\approx 3.76\times10^{-10}\,\mathrm{m} = 3.76\,\text{\AA}.
\end{align}

铜为单价金属，每个原子贡献 1 个自由电子。电子数密度
\begin{equation}
n = \frac{4}{a^3} \approx \frac{4}{5.32\times10^{-29}} \approx 7.52\times10^{28}\,\mathrm{m^{-3}}.
\end{equation}

\textbf{Step 2: 费米波矢与费米能}

自由电子模型中，考虑自旋简并：
\begin{equation}
n = \frac{k_F^3}{3\pi^2}
\quad\Longrightarrow\quad
k_F = (3\pi^2 n)^{1/3} = \Bigl(\frac{12\pi^2}{a^3}\Bigr)^{1/3} = \frac{(12\pi^2)^{1/3}}{a}.
\end{equation}
数值上 $(12\pi^2)^{1/3}\approx(118.4)^{1/3}\approx 4.91$，故
\begin{equation}
k_F \approx \frac{4.91}{3.76\times10^{-10}} \approx 1.31\times10^{10}\,\mathrm{m^{-1}}.
\end{equation}

绝对零度下费米能
\begin{equation}
E_F = \frac{\hbar^2k_F^2}{2m}
= \frac{(1.055\times10^{-34})^2\times(1.31\times10^{10})^2}{2\times9.11\times10^{-31}}
\approx 1.05\times10^{-18}\,\mathrm{J}.
\end{equation}
换算为电子伏特（$1\,\mathrm{eV}=1.602\times10^{-19}\,\mathrm{J}$）：
\begin{equation}
\boxed{E_F \approx \frac{1.05\times10^{-18}}{1.602\times10^{-19}} \approx 6.5\,\mathrm{eV}.}
\end{equation}

\textbf{Step 3: 导带宽度的估算}

在自由电子模型下，导带宽度可粗略估计为第一布里渊区中从中心到边界最远点的能量差。FCC 的第一布里渊区为截角八面体，最远点位于正方形面与六边形面交棱的顶点（$W$点）。

$W$ 点的坐标为 $(2\pi/a,\pi/a,0)$（以倒格子基矢为单位，或等价地，在直角坐标中为 $(\pi/a,2\pi/a,0)$？需统一）。标准坐标下，$W$ 点到原点的距离满足
\begin{equation}
k_W^2 = \Bigl(\frac{\pi}{a}\Bigr)^2 + \Bigl(\frac{2\pi}{a}\Bigr)^2 = \frac{5\pi^2}{a^2},
\qquad
k_W = \frac{\sqrt{5}\,\pi}{a}.
\end{equation}

对应最大波矢的能量即导带顶部能量（自由电子近似）：
\begin{equation}
E_{\max} = \frac{\hbar^2k_W^2}{2m}
= \frac{\hbar^2}{2m}\cdot\frac{5\pi^2}{a^2}
= \frac{5\pi^2\hbar^2}{2ma^2}.
\end{equation}

导带底部在 $\Gamma$ 点，能量为 $0$（以导带底为能量零点）。故导带宽度
\begin{equation}
\Delta E = E_{\max} = \frac{5\pi^2\hbar^2}{2ma^2}.
\end{equation}

代入数值：
\begin{align}
\Delta E &= \frac{5\times\pi^2\times(1.055\times10^{-34})^2}{2\times9.11\times10^{-31}\times(3.76\times10^{-10})^2} \notag\\
&\approx 2.13\times10^{-18}\,\mathrm{J}
\approx \boxed{13.3\,\mathrm{eV}.}
\end{align}
\end{derivationbox}

\begin{formulabox}[答案汇总]
\begin{center}
\begin{tabular}{ll}
\toprule
\textbf{问题} & \textbf{答案} \\
\midrule
(1) 费米能 & $E_F \approx 6.5\,\mathrm{eV}$ \\
(2) 导带宽度 & $\Delta E \approx 13.3\,\mathrm{eV}$ \\
\bottomrule
\end{tabular}
\end{center}
\end{formulabox}

\begin{insightbox}[物理图像与讨论]
\begin{itemize}
    \item \textbf{费米能约为导带宽度的一半}（$E_F/\Delta E\approx 0.5$）。这说明单价金属的传导电子仅填充了约半个导带，导带中有大量空态，电子在外场下很容易获得能量跃迁到空态，从而导电——这正是金属良好导电性的能带论解释。
    \item \textbf{自由电子模型的适用性}：铜的 $3d$ 带与 $4s$ 带实际上有交叠，自由电子模型对 $s$ 电子的估算在数量级上是合理的，但严格的能带计算需要更复杂的方法。
    \item \textbf{数量级记忆}：典型金属的费米能在几个 eV 量级（$1\!\sim\!10\,\mathrm{eV}$），导带宽度在 $10\!\sim\!20\,\mathrm{eV}$ 量级。铜的 $E_F\approx 7\,\mathrm{eV}$（实验值约 $7.0\,\mathrm{eV}$），与自由电子模型估算的 $6.5\,\mathrm{eV}$ 相当接近。
\end{itemize}
\end{insightbox}
